\section{Introduction}
\label{sec:intro}

La coloscopie est l'examen de référence pour le dépistage du cancer colorectal.
Les systèmes d'aide à la détection (CAD) analysent des images fixes pour signaler un polype suspect et confirmer les repères anatomiques requis par les scores de qualité.
Les réseaux convolutifs conviennent au déploiement quasi temps réel, mais les résultats publiés restent difficiles à comparer entre architectures, régimes de pré-entraînement et définitions de tâches.

Les benchmarks publics mélangent détection binaire, dépistage trois classes agrégé et taxonomies multi-classes sans protocole ni code partagés.
Le pré-entraînement auto-supervisé peut réduire le coût d'annotation, mais son avantage sur le transfert ImageNet en endoscopie reste une question empirique.

Les architectures prédictives jointes (JEPA)~\cite{lecun2022path,assran2023ijepa} apprennent en prédisant des représentations latentes de vues masquées.
Nous adaptons un JEPA \emph{convolutif} aux frames coloscopiques avec des blocs SE sensibles au masque compatibles avec le masquage spatial.

\frbrief{Nous fournissons un \textbf{benchmark et une implémentation de référence reproductibles}---sans prétendre que le JEPA domaine bat ImageNet sur les données publiques Simula.
La contribution principale est le protocole, le code et un rapport honnête sur les tâches complétées T3 (dépistage) et T1 (HyperKvasir-23).}

\textbf{Contributions.}
(1)~\textbf{Méthode :} CNN-JEPA sur SE-ResNet-50 avec attention canal sensible au masque.
(2)~\textbf{Benchmark :} C0--C5 et efficacité en labels C7--C9 (T3), plus B2--B6 (T1), avec arrêt anticipé, métriques JSON et tableaux auto-générés.
(3)~\textbf{Étude empirique :} le transfert ImageNet surpasse le CNN-JEPA domaine sur les splits publics actuels ; nous documentons les tendances d'efficacité en labels et les ablations sans surestimer les gains.
